(1) 径向惯性体力为 \(\rho\Omega^2r\)，平衡方程为 \(\sigma_r'+(\sigma_r-\sigma_\theta)/r+\rho\Omega^2r=0\)。轴对称应变 \(\epsilon_r=u'\)、\(\epsilon_\theta=u/r\)，平面应力关系为 \(E\epsilon_r=\sigma_r-\nu\sigma_\theta\)、\(E\epsilon_\theta=\sigma_\theta-\nu\sigma_r\)。

(2) 径向方程的正则解为 \(\sigma_r=A-(3+\nu)\rho\Omega^2r^2/8\)、\(\sigma_\theta=A-(1+3\nu)\rho\Omega^2r^2/8\)，舍弃中心发散的 \(1/r^2\) 分支。令 \(\sigma_r(R)=0\) 得 \(A=(3+\nu)\rho\Omega^2R^2/8\)，所以 \(\sigma_r=(3+\nu)\rho\Omega^2(R^2-r^2)/8\)。

(3) 利用 \(u/r=(\sigma_\theta-\nu\sigma_r)/E\) 可得 \(u(r)=\rho\Omega^2r[(1-\nu)(3+\nu)R^2-(1-\nu^2)r^2]/(8E)\)。在自由外缘 \(u(R)=\rho\Omega^2R^3(1-\nu)/(4E)>0\)，单位是长度且符号符合离心拉伸。

(4) 对 \(\nu\ge0\)，应力最大值在中心，\(\sigma_{\max}=(3+\nu)\rho\Omega^2R^2/8\)。故模型下 \(\Omega\le\sqrt{8\sigma_y/[(3+\nu)\rho R^2]}\)。达到塑性、显著厚度应力或 \(u/R\) 不再足够小时，必须放弃这里的薄盘线弹性准则。
